%% ma: L12-B01-0041
%% lop: 12
%% chuong: 1
%% bai: 1
%% dang: 12.1.2
%% loai: TN4
%% muc: TH
%% dap_an: "C"
%% nguon: "Ảnh giáo viên gửi 05/11/2026 (M11): Kiểm tra giữa HK1 Toán 12 (2026-2027), THPT Hồ Thị Bi (TP.HCM), mã đề 1133, Phần I Câu 8"
%% kien_thuc: "Bài 1 (tính đơn điệu: đọc từ đồ thị)"
%% trang_thai: chinh-thuc
%% ghi_chu: "Đề gốc không kèm đáp án; đã giải độc lập. Hình dựng lại từ hàm số y=x^4-4x^2+2 (cực tiểu (±√2;-2), cực đại (0;2)) khớp các số liệu trên hình gốc. kho.py trung báo trùng với L12-B01-0002, 0003, 0034 vì số nằm trong hình; hàm số, hình và phương án khác nên không trùng."
\begin{ex}
Cho hàm số có đồ thị như hình vẽ bên.
\begin{center}
\begin{tikzpicture}[scale=.9]
  \begin{scope}
    \clip (-2.4,-2.9) rectangle (2.4,3.2);
    \draw[thick,domain=-2.2:2.2,samples=100,smooth] plot (\x,{\x*\x*\x*\x-4*\x*\x+2});
  \end{scope}
  \draw[truc] (-2.5,0)--(2.6,0) node[below]{$x$};
  \draw[truc] (0,-2.9)--(0,3.3) node[left]{$y$};
  \draw[phu] (-1.4142,0)--(-1.4142,-2)--(1.4142,-2)--(1.4142,0);
  \path (0,0) node[below left]{\footnotesize$O$}
        (-1.4142,0) node[above left]{\footnotesize$-\sqrt2$} (1.4142,0) node[above right]{\footnotesize$\sqrt2$}
        (0,2) node[above right]{\footnotesize$2$} (0,-2) node[below right]{\footnotesize$-2$};
\end{tikzpicture}
\end{center}
Hàm số đã cho đồng biến trên khoảng nào sau đây?
\choice{$(0;+\infty)$}{$(-\sqrt2;2)$}{\True $(-\sqrt2;0)$}{$(-2;2)$}
\loigiai{Từ đồ thị, hàm số có cực tiểu tại $x=-\sqrt2$ và $x=\sqrt2$, cực đại tại $x=0$.
Đồ thị đi lên (từ trái sang phải) trên $(-\sqrt2;0)$ và $(\sqrt2;+\infty)$.
Trong các khoảng đã cho chỉ có $(-\sqrt2;0)$ nằm trọn trong khoảng đồng biến.}
\end{ex}
